\section{结论}
本文围绕可信事件记录能够约束多少真实估计损伤这一问题整合已有结果。
精确实际预测反馈将当前隐藏偏差带入下一可信事件的条件概率；
结合历史协方差校准与有限峰值，可得到全允许因果策略的统计损伤上界，
合法单拍构造进一步支撑双预算匹配阶数。

反馈量化扩展揭示了保证的一项边界。在标量、固定有限目标时刻及一致静默更新下，
内部预测可在同一量化格中变化而保留全部已指定的可信事件/量化反馈历史，
其最坏零记录TV损伤为模型相关的幅值--精度匹配阶数。
这一结论同时依赖量化中心参与触发和审计记录受限，
不说明完整接收日志对所有策略都无区分能力，也不说明增强审计已经消除全部损伤。

本稿作为理论安全性能分析，不强制依赖某一产品已经部署该协议；
但未验证具体工业部署、不提供部分认证的成本优势。
精确反馈界的数值常数及有限样本认证仍有明显保守性，
量化器的无饱和格点也不能视作固定有限比特实现。
新增冻结标量验证通过了时序、静默更新、量化记录路径保持和目标损伤的独立一致性检查，
并展示连续审计信息对当前见证的区分作用。
验证只覆盖精度主导区域的预设见证；当前样本未分辨精确反馈下极小的事件变化，
不据此声称零TV、最坏值数值求解或实用在线认证。
论文整体的新颖性、理论价值和投稿成熟度仍需在完整稿件层面评估。
